\chapter*{Relation to Prior Work}
\addcontentsline{toc}{chapter}{Relation to Prior Work}

This monograph sits downstream of several prior works, verified to
several different degrees of confidence, and this note distinguishes
them explicitly rather than let a reader assume uniform rigor. Four
sources have been checked by reading primary text directly:
\emph{Conservation of Ambiguity}, Distinguishability Geometry,
Persistence Before Truth, and Observability Is Restorability. Two
further sources, \emph{Admissibility Before Truth} and \emph{The
Geometry of Correction}, are integrated in Chapter~\ref{ch:four-operations}
on the strength of exact formulas and
theorem statements transcribed directly from the source PDFs by the
author, which is stronger evidence than paraphrase but is not the same as
an independent reading of the primary file; every claim drawn from them is marked
\transcribed{} rather than \imported{} throughout, and this distinction
is maintained rather than allowed to blur after the first citation.

\textbf{A retraction, since partially corrected.} An earlier draft of
this monograph cited \emph{The Geometry of Reconstruction} as the source
of the observation that repair operators require higher-order
maintenance (\cref{prop:confidence-renderings}). No such text has been
located, and that citation remains retracted. \emph{The Geometry of
Witnesses}, initially flagged as possibly forthcoming or nonexistent, is
now confirmed by the author to exist and to have been read; it is not,
however, transcribed or independently verified against primary source
text, and what this
monograph reports from it (\S\ref{sec:witnesses}) rests on the author's
prose summary rather than exact formulas copied from source text, unlike
\emph{Admissibility Before Truth} and \emph{The Geometry of Correction}.
It is cited at this third, weaker tier accordingly, and the distinction
is maintained rather than allowed to blur after the first citation.

From \emph{Observability Is Restorability} this monograph imports the
restorable-distinction set $\restdomain$ and the basic notion of repair
cost, both checked against primary text in
Chapter~\ref{ch:four-operations} (\S\ref{sec:oir}). That chapter also
reports a finding that bears on this monograph's own repair-interference
definition (\cref{def:interference}): \emph{Conservation of Ambiguity}
attributes its derivative-form interference functional to Observability
Is Restorability, but Observability Is Restorability's own text defines
interference algebraically instead. This monograph's usage is correctly
attributed to \emph{Conservation of Ambiguity}, where the derivative form
is actually stated; the discrepancy between the two source texts is
recorded in \cref{rem:interference-inconsistency} as a fact about the
corpus, not a defect in this monograph's citation.

From \emph{Admissibility Before Truth}, transcribed rather than
independently read, this monograph imports the Domain Restriction
Theorem (\S\ref{sec:abt}) and connects it directly to this monograph's
own sufficiency apparatus (\cref{prop:sufficiency-presupposes-evaluability}).
This text should not be confused with \emph{The Admissibility Field},
the presumed foundational text of the wider Admissibility Program, which
remains wholly unseen; \emph{Admissibility Before Truth} is a distinct,
narrower paper that happens to share the program's name.

From \emph{The Geometry of Correction}, likewise transcribed rather than
independently read, this monograph imports the distinction between
self-referential and reality-anchored evaluation and the coupling
coefficient $\rho$ (\S\ref{sec:goc}), connecting corruption
(\cref{def:corrupting-loss}) to correction suppression
(\cref{prop:corruption-as-suppression}).

From \emph{Conservation of Ambiguity} this monograph imports the
Redistribution Theorem in full and without modification
(\cref{thm:redistribution}), together with its explicit corollary that
ambiguity conservation is only a degenerate limiting case of the more
general repair-budget conservation. This monograph does \emph{not}
claim to be a corollary of that theorem. The Redistribution Theorem
governs a single observer's competing distinctions at one moment; this
monograph's Composition Failure Theorem
(\cref{thm:composition-failure}) governs a single distinction's status
as it moves through a sequence of independently-budgeted stages over
time. Appendix~\ref{app:comparison} gives the side-by-side comparison
in full, including a worked demonstration that one result is not
derivable from the other without the additional budget-relay
machinery developed in Part~II.

Distinguishability Geometry and Persistence Before Truth are imported with
verified source text in Chapter~\ref{ch:four-operations}, which states
precisely how far this monograph's independently-built vocabulary --
renderings, discard sets, root recovery, relay revision -- corresponds to
their actual formal apparatus, and where the correspondence is only
partial or remains open. The Admissibility Program's own foundational
text, \emph{The Admissibility Field}, has still not been consulted
directly in any form; this monograph's admissible continuation set
(\cref{def:admissible-set}) is accordingly flagged in
\cref{sec:collisions} as an unresolved correspondence. Three further,
distinct formal notions bearing on admissibility are now on record
without being reconciled to it or to each other: Observability Is
Restorability's observer-relative admissibility manifold $\mathcal{A}_O$,
the activation-manifolds application paper's admissibility field
$A: M \to \mathbb{R}_{\ge 0}$, and \emph{Admissibility Before Truth}'s
domain restriction $\mathrm{dom}(\mathcal{T}) \subseteq \mathrm{Rest}(\mathcal{R})$.
Of the three, the last connects most directly to this monograph's own
apparatus, as \S\ref{sec:abt} shows.
